\setchapterpreamble[u]{\margintoc}
\chapter{Partial Derivatives and Gradients}
\labch{gradients}

\sources{Hass, Heil and Weir, \emph{Thomas' Calculus}, 14th ed.\ \cite{thomas2018calculus},
Chapter~14; MIT 18.02 \cite{ocw1801}.}

A loss depends on many parameters at once, so the derivative has to become a
vector. This chapter builds that vector and computes it for the handful of
expressions that account for most of what appears in practice.

\section{Partial derivatives}
\labsec{partials}

For \(f:\R^{n}\to\R\), the partial derivative \(\partial f/\partial x_i\) is the
ordinary derivative with respect to \(x_i\) with the other coordinates held
fixed. Nothing new is required to compute one; the novelty is in assembling
them.

\begin{definition}
The \emph{gradient} of \(f\) at \(\vect{x}\) is the column vector
\[
	\nabla f(\vect{x})
	=\begin{bmatrix}
		\partial f/\partial x_1\\ \vdots\\ \partial f/\partial x_n
	\end{bmatrix}.
\]
\end{definition}

The first-order approximation of \refsec{derivative-definition} becomes
\begin{equation}
	f(\vect{x}+\vect{h})=f(\vect{x})+\nabla f(\vect{x})\tp\vect{h}+o(\lVert\vect{h}\rVert),
	\labeq{first-order}
\end{equation}
so the gradient is the vector whose inner product with a displacement predicts
the change in \(f\).

\section{Direction of steepest ascent}
\labsec{steepest}

\begin{proposition}
Among unit vectors \(\vect{u}\), the directional derivative
\[
	\nabla f(\vect{x})\tp\vect{u}
\]
is largest when \(\vect{u}\) points along \(\nabla f(\vect{x})\), and its value
there is \(\lVert\nabla f(\vect{x})\rVert\).
\end{proposition}

\begin{proof}
By the Cauchy--Schwarz inequality,
\(|\nabla f\tp\vect{u}|\le\lVert\nabla f\rVert\lVert\vect{u}\rVert
=\lVert\nabla f\rVert\), with equality exactly when \(\vect{u}\) is parallel to
\(\nabla f\).
\end{proof}

\marginnote{Hence \(-\nabla f\) is the direction of steepest \emph{descent}. That
single sentence is the justification for every method in
\refch{gradient-methods}.}

A related fact is that the gradient is perpendicular to the level set through
\(\vect{x}\): moving along a level set does not change \(f\), so the inner
product in \refeq{first-order} must vanish for such displacements.

\section{Gradients of the expressions that matter}
\labsec{matrix-gradients}

Three formulas cover most cases encountered in this book. Throughout,
\(\vect{a}\) is a fixed vector and \(A\) a fixed matrix.

\begin{proposition}
\begin{align*}
	\nabla\bigl(\vect{a}\tp\vect{x}\bigr)&=\vect{a},\\
	\nabla\bigl(\vect{x}\tp A\vect{x}\bigr)&=(A+A\tp)\vect{x},\\
	\nabla\Bigl(\tfrac12\lVert A\vect{x}-\vect{b}\rVert^{2}\Bigr)
		&=A\tp(A\vect{x}-\vect{b}).
\end{align*}
In particular the middle formula is \(2A\vect{x}\) when \(A\) is symmetric.
\end{proposition}

\begin{proof}[Third formula]
Write \(f(\vect{x})=\tfrac12(A\vect{x}-\vect{b})\tp(A\vect{x}-\vect{b})\) and
expand at \(\vect{x}+\vect{h}\):
\[
	f(\vect{x}+\vect{h})=f(\vect{x})
	+(A\vect{x}-\vect{b})\tp A\vect{h}
	+\tfrac12\lVert A\vect{h}\rVert^{2}.
\]
The middle term is \(\bigl(A\tp(A\vect{x}-\vect{b})\bigr)\tp\vect{h}\) and the
last is \(O(\lVert\vect{h}\rVert^{2})\); comparing with \refeq{first-order}
identifies the gradient.
\end{proof}

\begin{remark}
Setting the third gradient to zero gives \(A\tp A\vect{x}=A\tp\vect{b}\) --- the
normal equations of \refeq{normal-equations}. The geometric derivation in
\refch{orthogonality} and the calculus derivation here produce the same
equations, as they must.
\end{remark}

\section{A worked parameter gradient}
\labsec{worked-gradient}

Take one logistic unit with squared loss on a single example \((\vect{x},y)\):
\[
	z=\vect{w}\tp\vect{x}+b,
	\qquad
	\hat y=\sigma(z),
	\qquad
	\ell=\tfrac12(\hat y-y)^{2}.
\]
Working outwards from the loss,
\[
	\frac{\partial\ell}{\partial\hat y}=\hat y-y,
	\qquad
	\frac{\partial\hat y}{\partial z}=\hat y(1-\hat y),
	\qquad
	\frac{\partial z}{\partial\vect{w}}=\vect{x},
\]
so that
\[
	\nabla_{\vect{w}}\ell=(\hat y-y)\,\hat y(1-\hat y)\,\vect{x},
	\qquad
	\frac{\partial\ell}{\partial b}=(\hat y-y)\,\hat y(1-\hat y).
\]

\marginnote{Note the factor \(\hat y(1-\hat y)\), which is near zero whenever the
unit is confidently right or confidently wrong. Squared loss on a saturating
unit therefore learns slowly from its worst mistakes --- one of the reasons
cross-entropy is preferred, as \refch{estimation} explains.}

\section{Convexity in one line}
\labsec{convexity}

A differentiable \(f\) is \emph{convex} when
\[
	f(\vect{y})\ge f(\vect{x})+\nabla f(\vect{x})\tp(\vect{y}-\vect{x})
	\qquad\text{for all }\vect{x},\vect{y},
\]
that is, when the first-order approximation never overestimates. For convex
\(f\) a point with \(\nabla f=\vect{0}\) is a global minimum, which is why
convex problems are the ones that can be declared solved. The losses of interest
in Part~IV are not convex in the parameters, and \refch{gradient-methods} says
what survives of the theory anyway.
